% page 573 of the facsimile (image p-572 of the scan)
% block 167.6 x 276.5 mm ; left margin 34.4 mm
\pgc{12.700mm}{0.000mm}{
\l{\marge[78.403mm]{98.265mm}{\ornamento{11.551mm}{ornaments/letters/letter-T.png}}}
\l{}
\l{}
\l{}
\l{}
\l{\cel{53}565}
\l{}
\l{\sou{swabro}. (\sou{navo)}. Fasko de fili kanaba, por lavar e vishar la}
\l{\cel{3}ponto di navo. - DER.}
\l{}
\l{\sou{swichar}. (trans.) (\sou{elektro}). Chanjar la direciono di la elektro-}
\l{\cel{3}fluo, per aparato qua posibligas agar lo sen modifiko di la}
\l{\cel{3}instaluro, di la filo-reto. - E.}
\l{}
\l{\sou{+swobo.}\cel{2}Ligno-stipito, kartono-stipito, provizita,an lua extre-}
\l{\cel{3}maji, ye tufo piratra de kotono - uzata por skrapar ed un-}
\l{\cel{3}cionar la pelo di homo od animalo domestika. - E.}
\l{}
\l{}
\l{}
\l{}
\l{}
\l{}
\l{}
\l{}
\l{}
\l{\sou{ta.}\cel{2}(\sou{gram.}) Abreviuro di \textquotedbl{}ita\textquotedbl{}.}
\l{}
\l{\sou{tabako}. I. (\sou{bot.)}\cel{2}Planto ek la familio \textquotedbl{}solanei\textquotedbl{}. L.\cel{1}\sou{nicotiana}.}
\l{\cel{3}- II. Folio de ta planto, preparita por snifleso, fumeso (+smo-}
\l{\cel{3}keso) o mastikeso. - DEFIRS.}
\l{}
\l{\sou{tabano}. (\sou{zool.}) Insekto diptera, qua pikas la animali tante pro-}
\l{\cel{3}funde ke lia sango ekfluas. - L.\cel{1}\sou{tabanus}. IL.}
\l{}
\l{\sou{tabelo}. I. Panelo de ligno, indutita (maxim-multa-kaze) ye farbo}
\l{\cel{3}nigra o obskur-verda, e sur qua onu skribas o trasas figuri}
\l{\cel{3}per kreto-krayono. - II. Panelo ek ligno, kupro, e c., kadro}
\l{\cel{3}sur qua tensesis telo - sur qua onu piktas.- III. Listo di}
\l{\cel{3}qua la elementi (+itemi) esas ordinita metodoze (kom ex.:}
\l{\cel{3}tabelo de logaritmi). - DEFIRS.}
\l{}
\l{\sou{tabernaklo}. I. (\sou{antique, che la Judi}). Tendo en qua esis enklo-}
\l{\cel{3}zita la \textquotedbl{}archo di kontrato\textquotedbl{} til ante kande la templo esis kons-}
\l{\cel{3}truktita. - II. (\sou{liturgio katol.)}\cel{2}Recevuyo en qua esas enklo-}
\l{\cel{3}zita la santa ciborio, super la tablo altara. - DEFIS.}
\l{}
\l{\sou{tabeto}. (\sou{patol.}) Morbo quan karakterizas precipue la ne-inter-}
\l{\cel{3}koordineso di la movi, kun reteno di la forteso di la muskuli.}
\l{}
\l{}
\l{\sou{tablo}. Areo plata ek ligno, petro, e c., subtenata da un o plura}
\l{\cel{3}pedi, adsur qua onu pozas la objekti (por manjar, skribar,}
\l{\cel{3}Laborar, ludar, e c.) - EFI.}
\l{}
\l{\sou{tablero}. FS. (en\cel{1}\sou{automobilo}). Parieto, maxim-multa-kaze tola,}
\l{\cel{3}qua separas la kapoto (o : spaco rezervita por la motoro)}
\l{\cel{3}de la parto rezervita por la pasajeri. (Normale\cel{2}la automobil-}
\l{\cel{3}ilisto muntas sur la tablero la organi acesora : la exhaustilo,}
\l{\cel{3}la pumpilo di lubrifiko, la klaxono, e c.) - FS.}
\l{}
\l{tabulo. Planketo de ligno, en biblioteko, armoro, etajero, qua}
\l{\cel{3}interseparas horizontale e superpoze la parti diversa di la}
\l{\cel{3}kontenajo. - FIS.}
}
